\chapter{Semantic Identity}\label{ch03:identity}

This chapter develops the first and most foundational of the kernel's
five mechanisms: the semantic-identity layer. Its subject is the
question that the introduction's first failure class showed the
program cannot afford to leave implicit---\emph{when do two mentions
name the same thing?}---and its method is to make the answer a matter
of typed structure rather than of string coincidence or reader
discipline. The chapter proceeds from the five-field structural
reference and the rationale for each field
(Section~\ref{ch03:fivefield}), through the definition-hash binding
that detects definitional drift (Section~\ref{ch03:defhash}), the
disjoint namespace catalog (Section~\ref{ch03:namespaces}), the
link-kind system for cross-ontology correspondence and the historical
incident that motivates its central refusal
(Section~\ref{ch03:linkkinds}), the injective-rendering invariant
(Section~\ref{ch03:injective}), and the exact-type equality rule
(Section~\ref{ch03:exacttype}), to worked examples of legal and
illegal usage (Section~\ref{ch03:worked}) and a closing statement of
what the layer does not claim (Section~\ref{ch03:boundary}).

Throughout, the exposition is grounded in the reference
implementation: the module \texttt{src/cas\_core\_v8/semantics.py},
which defines the reference type, the namespace enumeration, the
crosswalk type, and the exact-match predicate, and the module
\texttt{src/cas\_core\_v8/identity.py}, which supplies the
validation primitives the reference type applies to its fields.
Where the chapter states what the code does, it reports the code as
read; where it states why, it reconstructs the design argument from
the governing specification and paper; and where the sources are
silent on a rationale, the chapter says so rather than supplying
one.

\section{The five-field structural reference}\label{ch03:fivefield}

\begin{definition}[Semantic reference]\label{ch03:def-semref}
A \emph{semantic reference} is the structural tuple
\[
(\text{namespace},\ \text{local\_id},\ \text{version},\
\text{source\_id},\ \text{definition\_sha256}),
\]
where the namespace is drawn from the closed catalog of
Section~\ref{ch03:namespaces}, \texttt{local\_id} and
\texttt{source\_id} are validated ASCII machine identifiers,
\texttt{version} is a non-blank canonical text, and
\texttt{definition\_sha256} is the SHA-256 digest, in normalized
uppercase hexadecimal, of the exact defining text of the identified
element in the identified source. Identity of references is
structural equality of the full tuple; no proper subset of the
fields identifies.
\end{definition}

In the implementation, Definition~\ref{ch03:def-semref} is the
frozen dataclass \texttt{SemanticRef} with exactly these five
fields, declared with \texttt{frozen=True} and \texttt{slots=True}
and validated exhaustively in \texttt{\_\_post\_init\_\_}: the
namespace must be an exact \texttt{SemanticNamespace} member;
\texttt{local\_id}, \texttt{version}, and \texttt{source\_id} must
be non-blank, trimmed strings free of forbidden Unicode categories;
\texttt{local\_id} and \texttt{source\_id} must additionally pass
the ASCII-identifier profile; and \texttt{definition\_sha256} is
normalized through the strict digest validator. A value that fails
any check does not exist as a \texttt{SemanticRef}: construction
raises \texttt{SemanticValidationError}. There is no partially
valid reference, no reference with a pending or empty field, and no
mutation path after construction. The docstring's one-line summary
is accurate enough to quote in paraphrase: one versioned semantic
object, bound to its defining source bytes.

The remainder of this section walks the five fields individually,
because each is there to kill a specific drift mechanism from the
introduction's taxonomy, and the design is best understood
field-by-threat.

\subsection{Namespace: which family}\label{ch03:field-namespace}

The namespace field answers the token-conflation problem directly.
The string \texttt{"F5"} is not a reference; a reference says
\emph{which family's} \texttt{F5}, as a typed enumeration member,
not as a naming convention. Because the field is an enumeration and
the constructor requires \texttt{type(namespace) is
SemanticNamespace}, there is no way to hold a reference whose
family is a guess: strings that merely spell a namespace's name do
not pass, and neither do subclass or look-alike enum members. Two
references with different namespaces are unequal regardless of
every other field, which is the formal rendering of the canon fact
that the families are disjoint. The deeper effect is on tooling:
any join, index, or cache keyed on references inherits the
namespace distinction for free, so the silent cross-family merge
that untyped strings invite becomes structurally impossible at
every downstream site, not just at disciplined ones.

\subsection{Local identifier: which element}\label{ch03:field-localid}

The \texttt{local\_id} names the element within its family. It is
deliberately a \emph{machine} identifier, validated against the
ASCII profile of Section~\ref{ch03:validation}: non-blank, trimmed,
ASCII-only, NFKC-stable, and drawn from an explicit identifier
alphabet. The profile exists because local identifiers are exactly
the strings humans review in diffs, registries, and receipts, and
the review channel is where invisible-character and confusable
spoofing does its damage: an identifier containing a zero-width
character or a bidirectional override renders identically to its
honest twin while being byte-distinct. The profile excludes the
entire class by construction---control, format, and surrogate
categories are rejected, non-ASCII is rejected, and
NFKC-instability is rejected---so that two local identifiers that
look the same to a reviewer are the same bytes. Prose names, by
contrast, are not local identifiers and do not belong in this
field; the layer's contract is that unrestricted text lives in
separate fields (such as the crosswalk note of
Section~\ref{ch03:linkkinds}) and never participates in identity.

\subsection{Version: which edition}\label{ch03:field-version}

The version field carries the edition of the element's definition
within its source---the answer to version drift. Its validation is
lighter than the identifier profile: the implementation requires a
non-blank, trimmed string free of forbidden Unicode categories, but
does not force the ASCII identifier alphabet, so version
designations follow the source's own versioning conventions. The
governing sources do not further constrain the version field's
format, and this chapter accordingly claims nothing about
comparability, ordering, or interpretation of version strings: two
references with different versions are simply different references.
That refusal to interpret is itself the point. A version field with
built-in ordering semantics would tempt tooling to treat
``newer'' as ``interchangeable-with''; the structural rule that any
version difference breaks identity forces every cross-version
inference through an explicit, reviewable act instead.

\subsection{Source identifier: which corpus}\label{ch03:field-sourceid}

The \texttt{source\_id} names the corpus or document in which the
defining text lives, under the same ASCII identifier profile as
\texttt{local\_id}. Its role is provenance at the identity layer:
the same element name, at the same version string, cited from two
different sources, is two different references. This is the field
that makes corroboration distinguishable from repetition---two
analyses grounded in the same source are visibly so---and it is
also the hook by which references connect to the program's
content-addressed artifact discipline, since a source identifier
can be resolved, in registries and receipts, to hash-bound source
bytes. The sources governing this chapter do not fix a global
catalog of source identifiers; the field's validity is syntactic,
and the binding of identifiers to actual corpora is registry
business outside this layer, a division of labor the chapter
returns to in Section~\ref{ch03:boundary}.

\subsection{Definition digest: which exact words}\label{ch03:field-defhash}

The fifth field binds the reference to the SHA-256 digest of the
exact defining text of the element in the identified source. It is
validated as exactly sixty-four hexadecimal digits and normalized
to uppercase, so that presentation-case differences cannot make
equal digests unequal. Its rationale is important enough to carry
its own section, next.

\subsection{Why five fields, and not fewer or more}\label{ch03:whyfive}

The tuple's arity is worth defending directly, because each of the
obvious alternatives fails against a named drift mechanism from
the introduction's taxonomy.

\emph{One field (the bare name)} is the status quo ante and the
proven failure: it answers none of the drift mechanisms and is the
representation in which the program's incidents occurred.

\emph{Two fields (namespace and name)} kill token conflation and
nothing else: definitional, version, and source drift all remain
invisible, and the flattened form remains collidable without a
rendering discipline.

\emph{An opaque surrogate (a UUID or database key)} is superficially
attractive---globally unique, cheap to compare---but it identifies a
\emph{row}, not a \emph{commitment}. Nothing in a surrogate is
checkable against the world: two parties cannot verify, from the
identifier alone, that they mean the same element as defined by the
same bytes, and every question about family, edition, or source
requires a resolution service that becomes a new trust dependency
and a new drift locus. The five-field tuple is deliberately
self-describing: every field is either directly meaningful or
directly checkable by rehashing, with no resolver in the loop for
the identity question itself.

\emph{A single content hash of the definition} (dropping the naming
fields) fails in the other direction: it commits to bytes but
cannot say what the bytes are \emph{for}. Two sources quoting the
same definitional sentence would collapse into one reference;
elements whose defining texts coincide (short definitions do
collide textually) would be conflated; and the reference would be
unreadable in review without a reverse lookup. Names without hashes
drift; hashes without names are mute; the tuple carries both
because identity here must serve both machines and reviewers.

\emph{More fields} were available---a digest of the whole source
corpus, a canon-edition field, a human-readable gloss---and the
governing sources do not record a rationale for each exclusion; the
chapter notes only the structural pattern in what was included:
every field present participates in identity and is validated by a
named grade of the string discipline, whereas glosses and corpus
digests belong to registries and provenance records, which already
carry them. Widening identity widens every equality comparison in
the system; the tuple is the minimal set that answers the taxonomy,
and minimality here is a reviewability feature.

\section{Definition-hash binding and definitional drift}\label{ch03:defhash}

\subsection{Naming a concept as defined by particular bytes}\label{ch03:asbytes}

The first four fields make a reference precise about family,
element, edition, and corpus. The fifth makes it precise about
something subtler: the actual words. A reference does not merely
name a concept; it names a concept \emph{as defined by particular
bytes}, and it carries a commitment to those bytes that any party
can check by rehashing the defining text.

To see why the first four fields are not enough, consider what
they leave open. A version string is an assertion by the source's
editors that a text is a particular edition; it is not a function
of the text. Sources get re-released with silent corrections under
unchanged version labels; working copies diverge from published
copies; a definitional sentence is touched in one checkout and not
another; a transcription introduces a variant that nobody flags as
an edition. In each of these cases, two parties can hold
references agreeing on all four leading fields while reading
different defining texts---which is precisely the definitional
drift of the introduction: identical names in good faith, divergent
meanings in fact. The digest closes the gap because it is computed
\emph{from} the text rather than asserted \emph{about} it. If the
defining bytes differ anywhere, the digests differ, the references
are unequal, and the drift that was invisible as a latent ambiguity
becomes visible as a plain inequality between two records.

\subsection{Staleness as a feature}\label{ch03:staleness}

The binding has a consequence that looks like a cost and is
actually the mechanism working: when a canon definition is
corrected, every reference bound to the superseded text becomes
\emph{stale}---its digest no longer matches what the current source
says. Nothing breaks silently; rather, every artifact carrying a
stale reference is now mechanically enumerable, which is exactly
the recheck obligation that untyped representations make
undiscoverable. The program's canon maintenance already imposes the
obligation in prose (derived artifacts are to be rechecked when a
governing definition changes); the digest field is what turns that
prose obligation into a query. It is worth being precise about what
staleness means and does not mean: a stale reference is not
invalid---it truthfully names the element as defined by the old
bytes, which is often exactly what a historical artifact should
say---but it is visibly non-current, and any pipeline that requires
currency can enforce it by comparing digests against the live
source. The layer itself takes no position on whether stale
references are acceptable in a given context; it only guarantees
that the question is decidable.

\subsection{What the digest does not do}\label{ch03:defhash-limits}

Three non-claims, in the spirit of the monograph's discipline.
First, the digest binds text, not meaning: two byte-identical
definitions are one definition, but the layer has no opinion on
whether two byte-different definitions are semantically equivalent
paraphrases---that is canon governance, not bookkeeping. Second,
the digest does not by itself prove that the named text is really
the defining text of the named element in the named source; a
reference constructed with a digest of the wrong passage is
well-formed and wrong. Verifying the binding requires resolving the
source and rehashing, which is registry and review business.
Third, the field inherits the general limits of hashing from the
background chapter: it is evidence about content, never about
authority, and its strength is the strength of SHA-256's collision
resistance, which the kernel consumes as standard background
rather than re-arguing.

\section{The namespace catalog}\label{ch03:namespaces}

\subsection{The catalog and its disjointness}\label{ch03:catalog}

The governing specification names seven disjoint namespace
families: TKS Foundation, TKS World, CAS Drive, CAS Domain, CAS
Register, CAS State, and Noetic Topos object. The implementation's
enumeration, \texttt{SemanticNamespace}, declares nine members:

\begin{itemize}
\item \texttt{TKS\_FOUNDATION}, \texttt{TKS\_WORLD}, and
\texttt{TKS\_NOETIC} on the TKS side;
\item \texttt{CAS\_DRIVE}, \texttt{CAS\_DOMAIN},
\texttt{CAS\_REGISTER}, \texttt{CAS\_NOETIC}, and
\texttt{CAS\_STATE} on the CAS side; and
\item \texttt{NOETIC\_TOPOS\_OBJECT} for the categorical
formulation.
\end{itemize}

The two extra members relative to the specification's prose list,
\texttt{TKS\_NOETIC} and \texttt{CAS\_NOETIC}, give the Noetic
family \emph{separate namespaces on each side of the TKS/CAS
divide}. The governing sources do not state the rationale in so
many words, and the chapter reports the natural reading as a
reading: the TKS Noetics and the CAS Noetics stand in exactly the
kind of positional, historically loaded correspondence that the
layer exists to keep honest, so folding them into one namespace
would build the layer's own failure mode into its catalog. Under
the nine-member enumeration, relating a TKS Noetic to its CAS
counterpart requires a typed crosswalk like any other cross-family
relation. Readers comparing specification and implementation
should treat the enumeration as the operative catalog; the
specification's seven-family list is not contradicted by it, since
the specification's ``Noetic Topos object'' family remains its own
namespace and the Noetic split refines rather than reshuffles the
prose taxonomy.

Table~\ref{ch03:tab-namespaces} records the operative catalog with
the family each member carries. The descriptions are orientation,
not definition: what each family's elements \emph{are} is fixed by
the program's canon documents, outside this layer, and the table
deliberately says only enough to keep the members apart.

\begin{table}[htbp]
\centering
\caption{The operative namespace catalog
(\texttt{SemanticNamespace}). Descriptions are orientational; the
canon owns the definitions.}
\label{ch03:tab-namespaces}
\begin{tabular}{@{}ll@{}}
\toprule
Member & Family carried \\
\midrule
\texttt{TKS\_FOUNDATION} & TKS Foundation positions \\
\texttt{TKS\_WORLD} & TKS World axis elements \\
\texttt{TKS\_NOETIC} & Noetic family, TKS side \\
\texttt{CAS\_DRIVE} & CAS motivational Drives \\
\texttt{CAS\_DOMAIN} & CAS processing Domains \\
\texttt{CAS\_REGISTER} & CAS Register axis values \\
\texttt{CAS\_NOETIC} & Noetic family, CAS side \\
\texttt{CAS\_STATE} & CAS Cognitive States \\
\texttt{NOETIC\_TOPOS\_OBJECT} & objects of the categorical
formulation \\
\bottomrule
\end{tabular}
\end{table}

Disjointness needs no runtime machinery: it is a property of the
type. An enumeration member belongs to exactly one enumeration,
every reference carries exactly one member, and equality of
references requires equality of members. There is no subtyping
among namespaces, no namespace wildcard, and no reference that is
``in'' two families. The catalog is also \emph{closed}: callers
cannot mint namespaces at run time, because the enumeration is the
type, and extending it is a source-level, review-visible change to
the kernel. Closure is a deliberate cost. A research program
discovers new families, and when it does, the kernel must change
through its gate rather than accept an open-world namespace
string---the design-inversion thesis applied to vocabulary.

\subsection{Namespaces are not a hierarchy}\label{ch03:nohierarchy}

One further negative deserves emphasis because the program's
ontologies invite it: the catalog encodes no relations among
namespaces. CAS States are, in the program's own canon, generated
by crossing Noetics with Domains; nothing in the namespace layer
represents that fact, and \texttt{CAS\_STATE} is not a composite,
child, or product namespace of \texttt{CAS\_NOETIC} and
\texttt{CAS\_DOMAIN} as far as identity is concerned. This is
intentional flattening. Structural relations among families are
semantic content, owned by the canon and representable---where
needed---by explicit typed links or by future reviewed extensions;
baking them into identity would make the identity layer a theory
of the domain, which is exactly what it must not be. The layer's
job is to keep names apart, not to explain them.

\section{Link kinds: relating without identifying}\label{ch03:linkkinds}

\subsection{The crosswalk object}\label{ch03:crosswalk}

Disjoint namespaces solve conflation and immediately create a new
need: the program's cross-family correspondences are real and
analytically load-bearing, so the layer must offer a way to state
them---one that makes stating them a deliberate, reviewable act.
The implementation's carrier is the frozen dataclass
\texttt{SemanticCrosswalk}, with fields: a validated ASCII
\texttt{crosswalk\_id}; \texttt{source} and \texttt{target}, each
an exact \texttt{SemanticRef}; a \texttt{relation} drawn from the
enumeration \texttt{CrosswalkRelation}; a
\texttt{source\_artifact\_sha256} digest identifying the artifact
that grounds the asserted relation; a mandatory non-blank
\texttt{note} of canonical text; and a research \texttt{boundary}
that the constructor requires to be exactly the frozen Ring-4 open
value---admission \textsc{open}, product \textsc{not licensed},
disclosure false---via the exact-match predicate
\texttt{is\_exact\_ring4\_open}.

Three of these fields enforce, at construction time, that a
crosswalk is an evidence-bearing assertion rather than a bare
arrow. The digest field means every crosswalk names the bytes of
the artifact on whose authority the correspondence is asserted:
there is no crosswalk ``because everyone knows.'' The mandatory
note means every crosswalk carries a human-readable justification;
an empty or whitespace note is a construction error, so the class
of undocumented links is empty by type. And the pinned Ring-4
boundary means no crosswalk can present itself as licensed or
disclosable: whatever a correspondence is, in this kernel it is
research material, and the object says so in a machine-checkable
field that the constructor refuses to let vary.

\subsection{The relation kinds}\label{ch03:kinds}

The enumeration \texttt{CrosswalkRelation} has four members:
\texttt{IDENTITY}, \texttt{TRANSLATION}, \texttt{ALIGNMENT}, and
\texttt{HYPOTHESIS}. The governing sources specify their intended
use at the level that matters for safety: a positional or
historical correspondence is represented as \textsc{translation},
\textsc{alignment}, or \textsc{hypothesis}, and is \emph{never}
silently represented as \textsc{identity}. The sources do not
supply formal definitions distinguishing the three non-identity
kinds from one another, and the chapter therefore offers only
their natural readings, marked as exposition: a
\emph{translation} link records a documented rendering of an
element of one system in another system's terms; an
\emph{alignment} records a structural or positional correspondence
between systems without asserting that either renders the other;
and a \emph{hypothesis} records a conjectured relation awaiting
review. What the layer enforces is not this gloss but the
machine-checkable part: the kind is a typed field, it is part of
the record, and no kind is coercible into another or into
identity. A consumer that treats a \texttt{HYPOTHESIS} link as if
it were a \texttt{TRANSLATION} is misusing a visibly typed record
---an auditable error, not a silent one.

\begin{ruleenv}[No silent identity]\label{ch03:rule-noid}
Cross-namespace identity is unrepresentable. The constructor of
\texttt{SemanticCrosswalk} permits the relation
\texttt{IDENTITY} only when the exact-match predicate
\texttt{semantic\_refs\_match} holds between source and target---
that is, only when the two endpoints are the very same versioned
reference, field for field. Since references in different
namespaces never match exactly, every cross-namespace crosswalk
must carry a non-identity kind, and the strongest cross-family
statement the type system can express is an explicitly typed,
noted, source-backed, Ring-4 link.
\end{ruleenv}

The implementation's error message for a violation is worth
reporting because it shows the historical failure written directly
into the code's contract: an \texttt{IDENTITY} crosswalk whose
endpoints differ raises a validation error stating that identity
is permitted only for the exact same versioned reference, and
naming the two canonical temptations---TKS Foundations related to
CAS Drives, and TKS Worlds related to CAS Domains---as pairs that
must use a non-identity relation. The degenerate self-link that
remains legal (an \texttt{IDENTITY} crosswalk from a reference to
itself) asserts nothing beyond what structural equality already
gives, and its admission keeps the enumeration honest: identity
exists as a kind, but the only thing it can ever relate is a
thing to itself.

\subsection{The crosswalk as a review artifact}\label{ch03:xw-review}

Read the crosswalk against the provenance
background of the survey chapter: the object is best
understood as a provenance record specialized to one assertion
kind. In provenance vocabulary, a crosswalk is an entity asserting
a relation between two other entities (the endpoint references),
derived from a grounding entity (the artifact named by
\texttt{source\_artifact\_sha256}), with the asserter's
justification carried in the note. What the kernel adds to the
generic pattern is fail-closed completeness: the generic
anti-patterns---derivation without a grounding entity, assertion
without justification, narrative provenance---are each excluded at
construction rather than discouraged in guidance.

This shape makes the crosswalk reviewable in a specific,
checklist-like way, and spelling the checklist out shows what the
type does and does not automate. A reviewer confronted with a
crosswalk record can ask: (i) do the endpoint references resolve
---do their sources exist and do their definition digests match
the cited passages? (ii) does the grounding artifact, fetched by
its digest, actually assert or evidence the claimed
correspondence? (iii) is the relation kind the strongest the
grounding artifact licenses, and in particular, would the
grounding artifact license more than \texttt{HYPOTHESIS}?
(iv) does the note fairly summarize the grounds? Questions (i)
and the digest half of (ii) are mechanical; questions (iii) and
(iv) are judgment. The type's contribution is that every judgment
question arrives with its materials attached and its scope
bounded---one relation, one grounds, one note---rather than
diffused through a document. The layer does not review; it makes
review tractable, which is the same division of labor the kernel
draws everywhere between structural consistency and human
authority.

A related operational point: because crosswalks are immutable
values, a correspondence that is later revised---upgraded from
\texttt{HYPOTHESIS} to \texttt{TRANSLATION} after review, or
withdrawn---is represented by a \emph{new} crosswalk (and, where
registries track it, a status change at the registry layer), not
by mutation of the old record. The old assertion remains what it
was, permanently citable as what was once claimed on what grounds.
The governing sources do not specify a crosswalk versioning or
supersession scheme beyond this immutability, and the chapter
accordingly leaves lifecycle policy to the registry chapters,
noting only that immutability is what makes any such policy
auditable.

\begin{finding}[The motivating incident]\label{ch03:finding-incident}
The rule's justification is historical before it is theoretical.
A prior work product in the program conflated positionally
corresponding Drive and Foundation names---treating the
correspondence as paraphrase rather than as relabel---and was
voided in review. The introduction analyzed why the artifacts of
the time could not resist the error: the names were untyped
strings, and no representation existed in which the correspondence
could be stated \emph{without} being stated as identity. The
crosswalk type is the direct answer. Under it, the same analytical
move the incident attempted is expressible---as an
\texttt{ALIGNMENT} or \texttt{TRANSLATION} with a note and a
source artifact---and the move the incident actually made is not
expressible at all.
\end{finding}

\begin{threat}[Namespace conflation]\label{ch03:threat-conflation}
An agent or tool treats a positional correspondence between
families as identity and propagates entailments across the
boundary. \emph{Blocked by} Definition~\ref{ch03:def-semref} and
Rule~\ref{ch03:rule-noid}: distinct namespaces make the endpoint
references unequal, and the crosswalk constructor refuses the
identity kind for unequal endpoints, so the conflation has no
representation; the nearest representable statement is a typed
non-identity link carrying its own note, source digest, and
research boundary.
\end{threat}

\section{The injective-rendering invariant}\label{ch03:injective}

\subsection{The problem with flattening}\label{ch03:flattening}

Structured references must sometimes travel as strings: as keys in
stores, as identifiers in logs and receipts, as line items in
review documents. Flattening is where a structurally sound design
can quietly reintroduce collision. The cautionary construction is
the naive delimiter join: render the tuple as
\texttt{namespace:local\_id:version:source\_id:digest} and two
distinct tuples collide whenever a component contains the
delimiter---the pair
(\texttt{a}, \texttt{b:c}) and the pair (\texttt{a:b}, \texttt{c})
flatten identically. An attacker who controls any component can
manufacture such collisions deliberately, aliasing one reference's
records to another's; an unlucky honest value can do it by
accident. Either way, every property the tuple design earned is
forfeited at the string boundary.

\begin{invariant}[Injective rendering]\label{ch03:inv-injective}
The rendered string form of a semantic reference is injective on
structural tuples: two distinct tuples render to distinct strings,
and delimiter-containing components cannot cause the rendering of
one tuple to equal the rendering of another.
\end{invariant}

\subsection{The design argument}\label{ch03:design-argument}

The implementation renders a reference through the property
\texttt{qualified\_id}, which serializes the five components, in
fixed order, as a JSON array with compact separators and
non-ASCII preserved, and prefixes the result with the versioned
scheme tag \texttt{semantic-ref.v1:}. The injectivity argument has
three steps, each carried by a specific feature of this
construction.

\begin{proposition}[Injectivity of the rendering]\label{ch03:prop-injective}
The map from valid \texttt{SemanticRef} tuples to their
\texttt{qualified\_id} strings is injective.
\end{proposition}

\begin{designargument}
First, \emph{escaping}: JSON string encoding escapes every
delimiter-significant character---quotation marks, backslashes,
and the structural commas and brackets appear inside components
only in escaped or quoted positions---so a component's content can
never terminate its own string early or open a neighboring one.
This is the length-prefixing role played generically: an encoding
in which component boundaries are unforgeable from component
content. Second, \emph{fixed arity and order}: the rendering is
always a five-element array in declared field order, so no
reading of the string can reassign bytes across field boundaries
or between differently-shaped tuples. Third, \emph{unique
parseability}: JSON array syntax parses unambiguously, so the
composite admits exactly one decomposition into components; two
tuples differing in any field therefore differ at the first
divergent component of the unique parse, and their renderings
differ as strings. The versioned prefix adds a fourth, forward-
looking guarantee: any future rendering scheme carries a
different tag, so strings from different schemes can never be
confused as one another. The property is enforced by construction
and is regression-tested with delimiter-bearing component
fixtures---components containing quotation marks, backslashes,
colons, commas, and bracket characters---whose pairwise renderings
the suite checks for distinctness.
\end{designargument}

\subsection{The design space: escaping, length-prefixing, hashing}\label{ch03:renderspace}

Three standard techniques achieve injective flattening, and the
choice among them is instructive. \emph{Length-prefixing} encodes
each component as its byte length followed by its bytes; boundaries
are unforgeable because they are computed, not delimited. It is the
technique of canonical binary formats, maximally simple to prove
correct, and minimally readable: a length-prefixed identifier is
opaque to a human reviewer, which matters in a layer whose strings
appear in diffs and review documents. \emph{Escaping} keeps the
composite readable---components appear as themselves, delimiters
are visible---at the price that correctness now depends on the
escaping rules covering every delimiter-significant character, in
every implementation that ever parses the form. Hand-rolled escape
schemes are historically where this goes wrong; adopting a
ubiquitous, rigorously specified escape discipline (here, JSON
string encoding) buys the readability of escaping while
outsourcing the hard part to a format with decades of
cross-implementation hardening and existing conformance suites.
\emph{Hashing the tuple} (rendering as a digest of a canonical
encoding) is injective up to collision resistance and produces
fixed-length names, but is doubly opaque---neither the fields nor
their boundaries survive---and forces every consumer through a
resolution step; it is the right shape for storage keys, not for
review-facing identifiers.

The implementation's choice---JSON-escaped, fixed-arity,
version-prefixed---sits deliberately at the readable end: a
\texttt{qualified\_id} can be eyeballed in a receipt, its fields
recognized, its scheme tag checked, and, when needed, parsed back
losslessly with a stock parser. The versioned prefix is the hedge
that keeps the choice revisable: a future scheme (length-prefixed,
hashed, or otherwise) would carry a different tag, and
Proposition~\ref{ch03:prop-injective}'s cross-scheme clause keeps
the generations from colliding. The governing sources do not
record a comparison of alternatives; the analysis above is the
chapter's reconstruction of the standard trade, offered so that a
reviewer can check the choice against its space rather than take
the space on faith.

Two remarks temper the proposition properly. The argument is a
design argument over the encoding, not a mechanized proof; the
cited fixtures exercise the adversarial component classes the
design anticipates, and the general claim inherits the usual
fixture-bound limits stated in the introduction. And injectivity
is claimed for the rendering of \emph{valid} references---the
validation layer has already excluded control characters and
unpaired surrogates from components, which removes the classes of
input on which JSON string encoding across implementations is
historically least uniform.

\begin{threat}[Identifier collision]\label{ch03:threat-collision}
Crafted components cause two distinct references to render to one
identifier, aliasing one reference's records, receipts, or
registry rows to another's. \emph{Blocked by}
Invariant~\ref{ch03:inv-injective} via
Proposition~\ref{ch03:prop-injective}: component content cannot
forge component boundaries, so the aliasing string does not
exist.
\end{threat}

\section{The exact-type equality rule}\label{ch03:exacttype}

\subsection{Virtual equality as an attack surface}\label{ch03:virtualeq}

In an object-oriented host language, equality is a method, and
methods can be overridden. This ordinarily benign fact becomes an
attack surface the moment equality carries security or evidence
weight. Consider a subclass of the reference type that overrides
equality and hashing: instances can report themselves equal to any
honest reference---matching it in every comparison, colliding with
it in every hash-keyed store---while carrying arbitrary divergent
field values and arbitrary overridden behavior. Every downstream
check that asks ``is this the reference I expected?'' through the
virtual equality protocol receives the answer the impostor chooses.
The same construction threatens boundary objects: a subclass of
the research-boundary type whose accessors report \textsc{open}
and \textsc{not licensed} while its stored state says otherwise
would pass any check that trusts the protocol. Nothing about this
requires exotic machinery; it is a dozen lines in any dynamic
language, and it defeats naive structural checking completely.

\begin{ruleenv}[Exact-type identity]\label{ch03:rule-exacttype}
Equality-overriding subclasses are not admitted as identity-bearing
values. Identity-bearing types reject subclass instances at
construction, and every load-bearing comparison of identity or
boundary values is performed by exact-type predicates that check
concrete runtime types and primitive fields directly, without
dispatching through overridable equality.
\end{ruleenv}

\subsection{Both halves in the implementation}\label{ch03:bothhalves}

The rule has two halves, and the implementation carries both. The
first half closes the front door: \texttt{SemanticRef},
\texttt{SemanticCrosswalk}, and the research-boundary type each
begin \texttt{\_\_post\_init\_\_} by checking
\texttt{type(self)} against the exact class and raising if a
subclass is being constructed; the artifact-identity type in
\texttt{identity.py} does the same. One cannot obtain an instance
of a subclass of any identity-bearing kernel type through its own
constructor.

The second half assumes the front door has failed---an impostor
object arrived from somewhere---and makes comparisons immune
anyway. The predicate \texttt{semantic\_refs\_match} compares two
candidate references by checking, in order: that each operand's
concrete type \texttt{is} exactly \texttt{SemanticRef}; that each
namespace's concrete type is exactly the namespace enumeration,
comparing their primitive values; and that each remaining field
is exactly a \texttt{str} on both sides and equal as such. At no
point does the predicate invoke \texttt{==} on the composite
objects; it never dispatches through anything an impostor could
override. The module applies the same pattern at its other
identity boundary (\texttt{artifact\_identities\_match} for byte
identities, field by field with type checks), and the crosswalk
constructor uses the exact predicate---not virtual equality---when
enforcing Rule~\ref{ch03:rule-noid}, so the no-silent-identity
rule cannot be satisfied by an object that merely claims equality.

\subsection{The discipline beyond references}\label{ch03:beyond}

The exact-type discipline is not local to the semantic layer; it
is the kernel's uniform equality posture, and seeing its second
instantiation clarifies the pattern. The byte-identity type
\texttt{ArtifactIdentity} in \texttt{identity.py}---the carrier of
the raw-digest and normalized-digest pair that the background
chapter's two-identity discipline requires---applies the same two
halves: its constructor rejects subclass instances, and its
comparison predicate, \texttt{artifact\_identities\_match}, walks
every primitive field with an exact \texttt{type(...) is} check
on both operands before comparing values, never invoking virtual
equality. The predicate is verbose by design---seven fields, each
with two type checks and one comparison---and the verbosity is
the audit trail: a reviewer can confirm, line by line, that no
comparison in the function dispatches anywhere overridable.

The pattern generalizes to a rule of thumb the later chapters
will keep instantiating: \emph{wherever a comparison gates
anything, the comparison is a named, exact-type predicate owned by
the kernel}. The alternatives each fail a review question.
Virtual equality fails ``who controls this comparison?''---the
answer is whichever operand's class got there first. Structural
equality via serialization (compare the encodings) fails ``what
does inequality mean?''---an encoding difference localizes poorly
and imports the whole encoder into every equality question.
Field-by-field comparison without type checks fails the impostor
case. The named-predicate pattern answers all three: the kernel
controls the comparison, inequality localizes to a field, and
impostors fail at the type gate. Its cost, beyond verbosity, is
that adding a field to an identity-bearing type requires touching
its predicate---which the kernel's authors evidently accepted as a
feature, since a field silently excluded from comparison is
exactly the kind of drift the discipline exists to prevent; the
constructor-side validation and the predicate enumerate the same
field lists, and a divergence between them is the sort of defect
the focused suites' fixtures are positioned to catch.

\begin{threat}[Equality spoofing]\label{ch03:threat-spoofing}
An adversarial or buggy subclass overrides equality and hashing to
impersonate an honest reference or boundary object, passing
downstream identity checks while carrying divergent state.
\emph{Blocked by} Rule~\ref{ch03:rule-exacttype}: construction of
the subclass through kernel constructors raises; and even a
foreign-built impostor fails every load-bearing comparison, because
the exact-type predicates check concrete types and primitive
fields and never call the overridable protocol.
\end{threat}

\begin{remark}[Cost of the rule]\label{ch03:remark-cost}
The rule spends real flexibility: well-behaved subclassing---a
convenience wrapper, a memoizing variant---is rejected along with
the malicious kind, and the exact-type predicates are verbose
compared to \texttt{==}. The trade is the design inversion again:
the kernel prefers a closed, verifiable equality discipline over
an open, convenient one, because the open one's safety would rest
on global knowledge that no subclass anywhere overrides equality
dishonestly---an unreviewable premise in exactly the way the
introduction's laundering premise was.
\end{remark}

\section{The validation pipeline}\label{ch03:validation}

The reference type's guarantees bottom out in three validation
primitives from \texttt{identity.py}, shared with the rest of the
kernel; the layer is best understood with their exact contracts in
view.

\texttt{require\_canonical\_text} accepts only exact \texttt{str}
values and rejects any character in Unicode categories
\texttt{Cc} (controls), \texttt{Cf} (format characters, including
zero-width and bidirectional controls), or \texttt{Cs}
(surrogates). It is the floor under every textual field in the
layer: even free-text fields like the crosswalk note, which accept
arbitrary languages and scripts, exclude the invisible-character
classes wholesale.

\texttt{require\_ascii\_identifier} layers the machine-identifier
profile on that floor: the value must be non-blank and equal to
its own trim; ASCII-only; a fixed point of NFKC normalization
(excluding compatibility characters that normalize to something
else than themselves); and a full match of the identifier
alphabet \texttt{[A-Za-z0-9][A-Za-z0-9.\_:/@+\textbackslash-]*}
---an alphanumeric head followed by alphanumerics and a small
punctuation set. This is the profile applied to
\texttt{local\_id}, \texttt{source\_id}, and
\texttt{crosswalk\_id}.

\texttt{require\_sha256} accepts only exact \texttt{str} values
matching sixty-four hexadecimal digits after uppercasing, and
returns the uppercase normalization; the stored digest field is
therefore case-canonical, and the presentation-discipline pitfall
of the background chapter (case-mismatched hexadecimal comparing
unequal) is closed at intake.

The division of labor among the three primitives reflects the
layer's string philosophy: three grades of string---digest,
machine identifier, canonical prose---each with a named validator,
each field of each type assigned exactly one grade, and no
load-bearing comparison ever performed on a string of the prose
grade.

\section{How the layer composes with the rest of the kernel}\label{ch03:composition}

Before the worked examples, a short tour of the layer's consumers,
since several of its design choices only make full sense in
composition.

\emph{Canonical encoding.} References and crosswalks are evidence
values, and both types expose the single declared serialization
route the encoding chapter requires: a \texttt{to\_dict} view whose
fields are exactly the validated primitives (plus, for references,
the \texttt{qualified\_id} rendering). The namespace and relation
enumerations pass through the encoder's enum tagging, so their
encodings are type-distinct from the plain strings that spell
their values---a reference whose namespace is the enumeration
member and a hand-built map whose namespace is the string
\texttt{"CAS\_DRIVE"} cannot hash equal. The validation grades of
Section~\ref{ch03:validation} also pre-satisfy the encoder's text
rules: by the time a reference exists, its strings already contain
no control or format characters, so a valid reference is always
encodable and the encoder never becomes a second, later point of
identity failure.

\emph{Receipts and registries.} Execution receipts and registry
records carry references in their identity fields, which is where
the layer's review-facing choices pay: a receipt naming its
semantic profile by \texttt{qualified\_id} is auditable by eye and
parseable by tooling, and a registry row keyed by rendered
reference inherits Invariant~\ref{ch03:inv-injective} as a
no-aliasing guarantee for its keying. Conversely, the layer leans
on the registries for everything it deliberately does not do:
resolving \texttt{source\_id} values to hash-bound corpora,
tracking crosswalk supersession, and enumerating stale references
after a canon change are registry operations over the immutable
values this layer mints.

\emph{Boundaries and the runtime gate.} Every crosswalk carries
the exact Ring-4 boundary, and the runtime decision model treats
identity-layer values as it treats all structural evidence:
sufficient to support nonclaim comparisons, never sufficient to
support claim-bearing intent. The layer thus supplies the
vocabulary for the kernel's refusals as well as its assertions---
when the gate denies a claim-bearing action, the record of the
denial names, in this layer's terms, exactly which semantic
objects were involved.

\section{Worked examples}\label{ch03:worked}

The following examples are expository constructions in the
implementation's vocabulary; digests are abbreviated with ellipses
for readability and would be full sixty-four-digit values in real
records.

\begin{example}[Legal: two references and a typed correspondence]\label{ch03:ex-legal}
An analyst needs to relate the fifth TKS Foundation, as defined in
a hash-bound edition of the TKS source corpus, to the CAS Drive
F5, as defined in a hash-bound edition of the CAS corpus. Two
references are constructed:

\texttt{SemanticRef(namespace=TKS\_FOUNDATION, local\_id="F5",}
\texttt{~~version="corpus-1.0", source\_id="tks-handbook",}
\texttt{~~definition\_sha256="93A1...C4")}

\texttt{SemanticRef(namespace=CAS\_DRIVE, local\_id="F5",}
\texttt{~~version="v7.7", source\_id="cas-canon",}
\texttt{~~definition\_sha256="0FE2...9B")}

The equal \texttt{local\_id} strings are harmless: the references
are unequal in namespace, source, and digest, and no kernel
operation will merge them. The correspondence is then stated as:

\texttt{SemanticCrosswalk(crosswalk\_id="xw-f5-drive5-0001",}
\texttt{~~source=<the TKS ref>, target=<the CAS ref>,}
\texttt{~~relation=ALIGNMENT,}
\texttt{~~source\_artifact\_sha256="7D10...EA",}
\texttt{~~note="Positional correspondence per the translation}
\texttt{~~key; not a paraphrase.")}

The crosswalk constructs successfully: endpoints are exact
references, the kind is non-identity, the note is non-blank, the
grounding artifact is named by digest, and the default boundary is
the exact Ring-4 open value. The record asserts precisely what the
canon licenses---a documented positional alignment---and nothing
stronger.
\end{example}

\begin{example}[Illegal: the conflation, replayed]\label{ch03:ex-illegal}
The same analyst---or an automated agent summarizing the
analyst's work---attempts the historical error: the same two
references with \texttt{relation=IDENTITY}. Construction raises
\texttt{SemanticValidationError}; the message states that identity
is permitted only for the exact same versioned reference and
names the Foundation/Drive and World/Domain pairs as required
non-identity cases. There is no flag to pass, no permissive mode,
and no alternative constructor. The error is not a policy lookup
that could be misconfigured; it is the type refusing to hold the
value.
\end{example}

\begin{example}[Illegal: malformed fields]\label{ch03:ex-malformed}
Each of the following constructions fails with a validation
error naming the offending field: a namespace supplied as the
plain string \texttt{"CAS\_DRIVE"} rather than the enumeration
member (namespace must be a \texttt{SemanticNamespace}); a
\texttt{local\_id} of \texttt{"~F5"} or \texttt{"F5~"}
(surrounding whitespace), \texttt{""} (blank), or \texttt{"F\,5"}
with a zero-width space between the characters (forbidden format
category---the case that is invisible in review and exists
precisely to be caught here); a \texttt{local\_id} containing
non-ASCII letters or failing NFKC stability; a
\texttt{definition\_sha256} of the wrong length, with non-hex
characters, or absent. A lowercase but otherwise well-formed
digest, by contrast, is \emph{accepted and normalized} to
uppercase---case is presentation, not identity.
\end{example}

\begin{example}[Illegal: the impostor]\label{ch03:ex-impostor}
A subclass \texttt{class LooksLikeRef(SemanticRef)} overriding
\texttt{\_\_eq\_\_} to return true against any reference cannot be
constructed: the base class's \texttt{\_\_post\_init\_\_} raises
on the subclass type before any field is validated. An impostor
built without the constructor still fails every kernel comparison,
because \texttt{semantic\_refs\_match} tests
\texttt{type(left) is SemanticRef} before touching any field. The
spoof has no purchase at either surface.
\end{example}

\begin{example}[Legal but delimiter-bearing: rendering under stress]\label{ch03:ex-delimiter}
A version string is prose-grade text and may lawfully contain
JSON-significant characters---say
\texttt{version='1.0 "corrected", second printing'}. The
\texttt{qualified\_id} rendering escapes the quotation marks
within the JSON string encoding, and no choice of other
components in any other reference can produce the same rendered
string: the delimiter fixtures of
Proposition~\ref{ch03:prop-injective} exercise exactly this class.
The naive colon-join of Section~\ref{ch03:flattening} would have
been ambiguous here; the versioned JSON rendering is not.
\end{example}

\begin{example}[Definitional drift, detected end to end]\label{ch03:ex-drift}
A final worked scenario exercises the digest field's purpose across
time. In month one, an analyst binds a reference to a canon
element: namespace \texttt{CAS\_DRIVE}, \texttt{local\_id}
\texttt{"F2"}, version \texttt{"v7.7"}, source
\texttt{"cas-canon"}, and the digest \texttt{"51BC...D3"} of the
element's defining sentence as it stands in the analyst's
hash-bound checkout. The reference enters a dozen analysis records
and one crosswalk. In month three, the canon receives a
corrected printing: the defining sentence of the same element is
amended---one clause tightened---without a change to the version
designation, an event the program's history shows is realistic.
A colleague working from the corrected corpus constructs the
``same'' reference and obtains digest \texttt{"E907...2A"}.

Nothing fails, and that is the design: both references are valid,
each truthfully names the element as defined by particular bytes,
and they are \emph{unequal}. The inequality surfaces exactly where
it should: a replay comparison between the analyst's records and
the colleague's reports a reference mismatch localized to the
digest field; a registry query for records bound to the old digest
enumerates precisely the dozen artifacts owing a recheck; and the
crosswalk, whose endpoint no longer matches the live source, is
flagged stale by the same comparison. Under a four-field design,
every one of these records would have continued to match while
meaning different words; the month-three divergence would have
been discovered, if ever, as a subtle analytical disagreement with
no mechanical cause. The digest field converts it into bookkeeping
---which is the entire thesis of the layer in one scenario.
\end{example}

\section{What the identity layer does not claim}\label{ch03:boundary}

\begin{boundary}[Limits of semantic identity]\label{ch03:bd-limits}
The layer guarantees, to the strength of its design arguments and
its cited fixtures: that references are structural, disjointly
namespaced, definition-bound, injectively renderable, and
comparable only by exact type and field; and that cross-namespace
identity is unrepresentable. It does not and cannot guarantee:
that a reference's digest was computed over the correct defining
passage (a registry and review obligation); that any crosswalk's
asserted correspondence is true (the crosswalk records an
assertion and its grounds, not a fact); that the namespace catalog
is complete for the program's future needs (extension is a
reviewed kernel change); or that any use of the layer confers
semantic correctness, empirical validity, or authority of any
kind. Every value the layer produces carries, or is governed by,
the Ring-4 research boundary---admission \textsc{open}, product
\textsc{not licensed}, disclosure withheld---and nothing in this
chapter's mechanisms alters that boundary. The bounded receipts
that exercise this layer (reference identity, rendering
distinctness under delimiter fixtures, crosswalk rejection cases,
exact-type rejections) are historical, fixture-bound software
evidence in the sense fixed by the introduction, and support no
claim beyond tested behavior under those fixtures.
\end{boundary}

\begin{remark}[The adversary model, stated plainly]\label{ch03:remark-adversary}
A reader may reasonably ask who the adversary of this chapter's
threat statements is supposed to be, since a research kernel is
not an internet-facing service. The honest answer, consistent
with the program's history, is threefold, in ascending order of
intent. First, \emph{entropy}: honest tools and honest authors
making locally reasonable choices---a different serializer, a
convenient string join, a subclass for a quick experiment---whose
composition produces conflation and aliasing with no adversary
anywhere; most of the layer's rules earn their keep against this
adversary alone. Second, \emph{the program's own generative
tooling}: AI agents that produce code, records, and summaries at
scale, whose failure mode is not malice but fluent overreach---
constructing the plausible-looking link, the confident identity,
the tidy rendering---faster than review can police; fail-closed
construction converts this pressure into visible errors at the
point of generation. Third, \emph{a deliberate adversary}: the
spoofing and collision constructions of this chapter are written
as attacks because design against the deliberate case covers the
accidental ones, and because a kernel meant to underwrite future
claims must assume its inputs will one day be crafted, not merely
careless. The threat statements name the third adversary; the
first two are the ones the program has actually met, and they are
covered a fortiori.
\end{remark}

The layer's contribution to the kernel is therefore best stated in
the design-inversion vocabulary: it does not detect conflation, it
\emph{disarms} it, by making the conflating statement
inexpressible and the honest statement cheap. Every mechanism in
the chapters that follow---boundaries, encodings, compatibility
receipts, registries---consumes references and crosswalks in the
forms fixed here, which is why the layer comes first: it is the
vocabulary in which the rest of the kernel is prevented from
lying.
